\begin{tikzpicture}[font=\sffamily\scriptsize]
% ================= the bench line =================
\draw[line width=1.1pt, draw=black!70] (-4.8,-1.15) -- (11.2,-1.15);
% ================= the Pi =================
\node[board, minimum width=40mm, minimum height=20mm] (pi) at (-3.0,0) {};
\node[text=white, font=\sffamily\bfseries\small] at (-3.0,0.25) {Raspberry Pi 4};
\node[text=white, font=\tiny] at (-3.0,-0.20) {the logger};
\node[draw=black!70, fill=gray!55, minimum width=4mm, minimum height=5mm, inner sep=0pt]
  (piusb) at (-1.0,0.35) {};
% ================= the PPK2 =================
\node[draw=linecol, fill=hwcol, rounded corners=1.5mm, minimum width=38mm,
      minimum height=20mm] (ppk) at (1.9,0) {};
\node[font=\sffamily\bfseries\small] at (1.9,0.45) {nRF PPK2};
\node[font=\tiny, align=center] at (1.9,-0.05) {source mode\\3300 mV, about 1 A};
\node[draw=black!70, fill=gray!55, minimum width=4mm, minimum height=5mm, inner sep=0pt]
  (ppkusb) at (0.1,0.35) {};
\node[font=\tiny, anchor=west] at (0.3,-0.75) {VIN};
\node[font=\tiny, anchor=east] at (3.5,-0.75) {VOUT};
% USB from the PPK2 to the Pi
\draw[usbcable] (ppkusb.west) -- (piusb.east);
\node[anchor=south, font=\tiny, align=center, fill=white, inner sep=1.5pt] at (-0.5,1.12)
  {PPK2 USB to the Pi};
% ================= the supply =================
\node[draw=linecol, fill=extcol, rounded corners=1.5mm, minimum width=26mm,
      minimum height=14mm] (sup) at (1.9,2.3) {};
\node[font=\sffamily\bfseries\small] at (1.9,2.5) {bench 5 V};
\node[font=\tiny] at (1.9,2.1) {or a Pi USB port};
\draw[cable, black] (1.9,1.65) -- (1.9,1.0);
\node[anchor=west, font=\tiny] at (2.0,1.35) {5 V into VIN};
% ================= the DUT =================
\node[draw=linecol, fill=usrcol, rounded corners=1.5mm, minimum width=34mm,
      minimum height=18mm] (dut) at (7.4,0) {};
\node[font=\sffamily\bfseries\small] at (7.4,0.35) {ESP32 NodeMCU};
\node[font=\tiny, align=center] at (7.4,-0.15) {idle, BLE advertise,\\deep-sleep};
% two flying leads from VOUT to the DUT
\draw[cable, red!70!black] (3.85,-0.55) to[out=0, in=180] (5.65,0.25);
\draw[cable, black!70]     (3.85,-0.85) to[out=0, in=180] (5.65,-0.25);
\node[anchor=south, align=center, font=\tiny] at (4.72,0.42) {3V3 and GND};
% ================= the DUT USB, kept off the bench =================
\node[draw=black!70, fill=gray!55, minimum width=3mm, minimum height=4.5mm, inner sep=0pt]
  at (9.15,-0.35) {};
\draw[line width=1.2pt, draw=red!70!black] (9.55,-0.58) -- (10.01,-0.12);
\draw[line width=1.2pt, draw=red!70!black] (10.01,-0.58) -- (9.55,-0.12);
\node[anchor=west, align=left, font=\tiny, text=red!65!black] at (10.15,-0.35)
  {no cable here};
% the cable, coiled somewhere else
\draw[usbcable, black!45] (6.6,-1.75) circle (0.32);
\draw[usbcable, black!45] (6.95,-1.75) circle (0.32);
\node[anchor=west, align=left, font=\tiny] at (7.45,-1.75)
  {the DUT USB cable lives here, off the bench,\\not coiled beside the board};
% ================= the POW-BB =================
\node[draw=linecol, fill=notecol, rounded corners=1.5mm, minimum width=26mm,
      minimum height=11mm] (pow) at (-3.0,-1.9) {};
\node[font=\sffamily\bfseries\tiny] at (-3.0,-1.72) {SBC-POW-BB};
\node[font=\tiny, align=center] at (-3.0,-2.05) {labelled rails,\\P15 discipline};
% ================= the checklist =================
\node[draw=linecol, fill=white, rounded corners=2pt, align=left, font=\tiny,
      text width=52mm, inner sep=4pt, anchor=north west] at (-1.2,-1.45)
  {before the first capture\\
   1. set source mode and 3300 mV with the DUT off\\
   2. check ground continuity with a meter\\
   3. confirm the DUT USB cable is out of reach\\
   4. then connect the DUT and start the run};
\end{tikzpicture}
